\chapter*{ओपन लॉजिक प्रकल्पाविषयी}
\addcontentsline{toc}{chapter}{ओपन लॉजिक प्रकल्पाविषयी}

\textit{Open Logic Text} हे औपचारिक अधितर्कशास्त्र आणि
औपचारिक पद्धती यांवरील मुक्त-स्रोत, सहयोगाने लिहिलेले पाठ्यपुस्तक
आहे. ते मध्यम स्तरापासून सुरू होते (म्हणजे औपचारिक तर्कशास्त्राचा
प्राथमिक अभ्यासक्रम पूर्ण झाल्यानंतर). गणितेतर वाचकांना, विशेषतः
तत्त्वज्ञान आणि संगणकशास्त्राच्या विद्यार्थ्यांना, उद्देशून लिहिलेले
असले तरी ते काटेकोर आहे.

सध्या समाविष्ट असलेल्या काही विषयांचा ऊहापोह अजून अपूर्ण
असू शकतो आणि अनेक विभागांत मोठी सुधारणा आवश्यक आहे. भविष्यात
अधिक विषय समाविष्ट करण्याचा आमचा विचार आहे. शब्दसूची, पुढील
वाचनाची यादी, ऐतिहासिक नोंदी, चित्रे, अधिक चांगली स्पष्टीकरणे,
तत्त्वज्ञान, संगणकशास्त्र आणि गणित यांमध्ये निष्कर्षांचे महत्त्व
उलगडणारे विभाग, तसेच अधिक प्रश्न आणि उदाहरणे यांचीही भर
घालण्याचा आमचा विचार आहे. तुम्हाला चूक आढळली किंवा काही सूचना
असली, तर \href{https://github.com/OpenLogicProject/OpenLogic/wiki/Contributing}{कृपया प्रकल्पाच्या चमूला कळवा}.

हा प्रकल्प मुक्त-स्रोताच्या भावनेने चालतो. मजकूर विनामूल्य
उपलब्ध तर आहेच; त्याबरोबर त्याचा LaTeX स्रोत क्रिएटिव्ह कॉमन्स
अट्रिब्यूशन परवान्याखाली उपलब्ध करून दिला आहे. त्यामुळे योग्य
श्रेय दिल्यास कोणालाही आमचे काम डाउनलोड करणे, वापरणे, बदलणे,
त्याची मांडणी बदलणे, दुसऱ्या स्वरूपात रूपांतर करणे आणि पुन्हा
वितरित करणे यांचा अधिकार आहे. अधिक माहितीसाठी ओपन लॉजिक
प्रकल्पाचे \href{http://openlogicproject.org/}{openlogicproject.org} संकेतस्थळ पहा.
